\section{Results}
\label{sec:results}

\subsection{Main Results}
\label{sec:main-results}

Table~\ref{tab:main-results} reports classification performance. Several
findings stand out.

\begin{table*}[!ht]
    \caption{Classification metrics for seven controlled Llama retrieval
    configurations and two no-retrieval Jev comparators (direct and two-stage
    behavior-guided). FPR is the false-positive rate. Bold values are the best
    observed values, not a controlled cross-model comparison.}
    \label{tab:main-results}
    \centering
    \begin{tabular}{lccc|ccc|cc}
        \toprule
        \textbf{Method} & \multicolumn{3}{c}{\textbf{Malicious}} & \multicolumn{3}{c}{\textbf{Benign}} & \multicolumn{2}{c}{\textbf{Overall}} \\
        \cmidrule(lr){2-4}\cmidrule(lr){5-7}\cmidrule(lr){8-9}
        & \textbf{Precision}~$\uparrow$ & \textbf{Recall}~$\uparrow$ & \textbf{F1 Score}~$\uparrow$ & \textbf{Precision}~$\uparrow$ & \textbf{Recall}~$\uparrow$ & \textbf{F1 Score}~$\uparrow$ & \textbf{BAcc}~$\uparrow$ & \textbf{FPR}~$\downarrow$ \\
        \midrule
        No RAG           & \textbf{1.000} & 0.798 & 0.888 & 0.948 & \textbf{1.000} & 0.973 & 0.899 & \textbf{0.0\%} \\
        Dense            & 0.972 & \textbf{0.972} & \textbf{0.972} & \textbf{0.991} & 0.991 & \textbf{0.991} & \textbf{0.981} & 0.9\% \\
        BM25             & 0.964 & \textbf{0.972} & 0.968 & \textbf{0.991} & 0.988 & 0.989 & 0.980 & 1.2\% \\
        Hybrid           & 0.976 & 0.964 & 0.970 & 0.988 & 0.992 & 0.990 & 0.978 & 0.8\% \\
        Behavior         & \textbf{1.000} & 0.782 & 0.878 & 0.931 & \textbf{1.000} & 0.964 & 0.891 & \textbf{0.0\%} \\
        Contrastive      & 0.996 & 0.948 & 0.971 & 0.982 & 0.999 & 0.990 & 0.973 & \textbf{0.1\%} \\
        Beh.\ + Contr.   & \textbf{1.000} & 0.788 & 0.881 & 0.933 & \textbf{1.000} & 0.966 & 0.894 & \textbf{0.0\%} \\
        Jev (direct)     & \textbf{1.000} & 0.922 & 0.960 & 0.974 & \textbf{1.000} & 0.987 & 0.961 & \textbf{0.0\%} \\
        Jev (beh.-guided)& 0.996 & 0.915 & 0.954 & 0.972 & 0.999 & 0.985 & 0.957 & 0.1\% \\
        \bottomrule
    \end{tabular}
\end{table*}

\paragraph{Finding 1: retrieval improves detection within the Llama suite.}
The zero-shot baseline achieves F1~0.888 with perfect precision but misses
roughly one in five malicious samples (recall 0.798). Dense retrieval lifts
malicious recall to 0.972 and F1 to 0.972 (+8.4 points) with a balanced
accuracy of 0.981; BM25 and hybrid essentially match this level (F1
0.968--0.970). Retrieval is thus useful for malicious-code detection, and the
gain comes almost entirely from exposing the model to similar \emph{code},
regardless of the similarity mechanism.

\paragraph{Direct Jev comparator.} Without retrieval, Jev obtains malicious F1
of 0.960 and balanced accuracy of 0.961, with no false positives. Its lower
malicious recall than the strongest retrieval configuration (0.922 versus
0.972) shows that source-only typed decisions remain sensitive to sparse or
ambiguous \texttt{setup.py} evidence. This comparison is indicative rather
than a controlled ablation because Jev is a different model family. Of the
1{,}020 evaluable records, nine unusually large malicious sources returned API
validation errors, yielding 0.991 coverage; failed calls are excluded from the
classification denominators and reported through coverage. The complete run
used 0.049 USD of OpenRouter Jev inference under the recorded provider/model
configuration.

\paragraph{Behavior guidance does not improve classification.} The two-stage
Jev variant reaches malicious F1 of 0.954 and balanced accuracy of 0.957,
slightly below direct Jev (0.960 and 0.961). It also introduces one false
positive while reducing malicious recall from 0.922 to 0.915. Thus, feeding
thresholded behavior indicators back into Jev does not improve this task under
the tested prompt and 0.5 threshold. The two-stage run costs 0.136 USD under
the recorded provider/model configuration and has the same 0.991 coverage as
direct Jev.

\paragraph{Finding 2: sparse retrieval is not worse than dense retrieval.}
BM25 (0.968) is within 0.4 F1 points of dense (0.972) at a fraction of the
engineering cost---an in-memory term index, no embedding model, no vector
store. Fusing both does not beat the better of the two, consistent with prior
RAG literature on short, keyword-heavy security queries. For this task,
lexical similarity over code tokens is a strong and cheap baseline that
future work should not skip.

\paragraph{Finding 3: contrastive evidence trades recall for the fewest
false alarms.} Contrastive retrieval, which shows the model both malicious
and benign reference groups, yields the best precision among RAG methods
(0.996) with FPR of 0.1\%---an order of magnitude below dense, BM25, and
hybrid (0.8--1.2\%)---at the cost of recalling slightly fewer malicious
samples (0.948 vs.\ 0.972). In a setting where triaging false alarms
dominates cost, contrastive retrieval is the preferred configuration.

\paragraph{Finding 4: static behavior summaries are lossy.} Replacing the
query with an AST-derived behavior summary, or adding it on top of
contrastive retrieval, actively hurts: F1 drops to 0.878/0.881, below the
no-RAG baseline in recall-adjusted terms, although all misclassifications
remain false negatives and FPR stays at zero. The summaries discard
identifier names, call structure and file-path information that the model
uses to recognize attacks such as typosquats; retrieving behavior-summary
neighbors does not compensate.

\subsection{Faithfulness of Explanations}
\label{sec:faithfulness}

Table~\ref{tab:faithfulness} reports the static-grounding proxy computed by
keyword-normalizing the model's free-form behavior claims to our static
vocabulary and checking each against \texttt{ast}-derived behavior flags.
Direct Jev returns a typed classification without behavior claims or source
citations, so its faithfulness statistics are not applicable and it is
omitted from the table. Behavior-guided
Jev does return typed behavior probabilities; its row is typed behavior
alignment rather than free-form explanation faithfulness and is therefore not
directly comparable with the Llama rows.
Unsupported-claim rates are high across the board (0.59--0.72), a caveat we
discuss in Section~\ref{sec:discussion}, but the ranking is informative: the
no-RAG baseline is the least grounded (0.717), all retrieval variants reduce
unsupported claims, and behavior-based queries yield the lowest rate (0.593)
with the best claim precision pair (behavior and contrastive: 0.387/0.421).
Retrieving explicit evidence therefore correlates with more grounded, less
extrapolated explanations, and behavior prompts in particular restrain the
model to security-relevant vocabulary.

\begin{table}[!ht]
    \caption{Automated static-grounding statistics for Llama model-claimed
    behaviors and typed behavior alignment for behavior-guided Jev. \texttt{Unsup.} =
    unsupported claimed behaviors / all claimed; \texttt{Beh.~P} = supported
    claims / claimed; \texttt{Beh.~R} = claimed-and-supported behaviors /
    statically observed behaviors.}
    \label{tab:faithfulness}
    \centering
    \setlength{\tabcolsep}{4pt}
    \begin{tabular}{lrrrr}
        \toprule
        \textbf{Method} & \textbf{Claims} & \textbf{Unsup.}~$\downarrow$ & \textbf{Beh.~P}~$\uparrow$ & \textbf{Beh.~R}~$\uparrow$ \\
        \midrule
        No RAG         & 427 & 0.717 & 0.282 & 0.427 \\
        Dense          & 560 & 0.657 & 0.312 & 0.521 \\
        BM25           & 596 & 0.685 & 0.294 & 0.517 \\
        Hybrid         & 568 & 0.665 & 0.314 & 0.525 \\
        Behavior       & 526 & \textbf{0.593} & 0.387 & \textbf{0.531} \\
        Contrastive    & 455 & 0.615 & \textbf{0.421} & 0.457 \\
        Beh.\ + Contr. & 697 & 0.694 & 0.289 & 0.543 \\
        Jev (beh.-guided)$^{\dagger}$ & 1{,}613 & 0.747 & 0.202 & 0.927 \\
        \bottomrule
    \end{tabular}
    \par\smallskip\footnotesize{$^{\dagger}$Typed behavior alignment at a 0.5
    probability threshold; Jev does not generate source-grounded explanations.}
\end{table}

\subsection{Error Analysis}
\label{sec:error-analysis}

\paragraph{Consistently missed samples.} Five malicious packages are
misclassified as benign by \emph{all seven Llama configurations}:
\texttt{parseweb-v3}, \texttt{kazer12-v2}, and \texttt{testdufou-v2} are
near-verbatim metadata clones of the legitimate \texttt{colorama} package
(same authors, description and project URL in \texttt{setup.py}) that differ
only by name; \texttt{pepequests-v0.0.1} is a minimal typosquat of
\texttt{requests} whose \texttt{setup.py} contains no executable behavior at
all; and \texttt{ucap-v3.6.1} hides its malicious setup behind innocuous
packaging boilerplate. One further typosquat, \texttt{simplyjson}, is missed
by six of the seven Llama configurations. These failures are structural: when the
\emph{specific file we analyze} carries no distinguishing signal, no
retrieval over the same file type can help, and name-level analysis (e.g.,
similarity to popular packages) would be required.

\paragraph{Recurring false positives.} Two benign packages are flagged as
malicious by three of the seven Llama configurations
(\texttt{authcaptureproxy-1.1.4},
\texttt{undetected-chromedriver-3.1.5.post4}),
both of which legitimately exercise capture-proxy and privilege-heavy browser
automation---code patterns orthographically similar to credential-theft
trojans in the knowledge base. Spot checks of the sampled false positives
suggest they arise from the model over-extending retrieved malicious
evidence, the failure mode that contrastive retrieval is designed to
suppress: indeed it flags neither of these two packages.

\paragraph{What the model sees.} Figure~\ref{fig:error-examples} contrasts
three real \texttt{setup.py} excerpts. The two false negatives contain
opposite forms of missing evidence: one impersonates a known benign package
through its metadata, while the other contains no executable installation
behavior. The false positive instead contains security-sensitive language in
a benign package declaration. Together, these examples show why source-only
signals can be insufficient even when retrieval improves aggregate recall.

\begin{figure*}[!ht]
    \centering
    \fcolorbox{black}{black!4}{%
    \begin{minipage}[t]{0.30\textwidth}
        \vspace{0.5ex}
        \textbf{A. Metadata clone}\\[-0.25ex]
        \scriptsize\textbf{Package:} \texttt{parseweb-v3}\\
        \scriptsize\textbf{Gold:} malicious \quad \textbf{Outcome:} FN (7/7)\\[0.75ex]
        \rule{\linewidth}{0.3pt}\\[-0.5ex]
        \footnotesize\ttfamily
        name='parseweb'\\
        author='Jonathan Hartley'\\
        description='Pure-python Colors'\\
        url='.../tartley/colorama'\\
        install\_requires=['requests', 'colorama']\\[0.75ex]
        \normalfont\scriptsize
        The analyzed file mirrors \texttt{colorama}; only the package name
        conflicts with the benign identity.
        \vspace{0.5ex}
    \end{minipage}}
    \hfill
    \fcolorbox{black}{black!4}{%
    \begin{minipage}[t]{0.30\textwidth}
        \vspace{0.5ex}
        \textbf{B. Behavior-empty typosquat}\\[-0.25ex]
        \scriptsize\textbf{Package:} \texttt{pepequests-v0.0.1}\\
        \scriptsize\textbf{Gold:} malicious \quad \textbf{Outcome:} FN (7/7)\\[0.75ex]
        \rule{\linewidth}{0.3pt}\\[-0.5ex]
        \footnotesize\ttfamily
        name = 'pepequests'\\
        version = '0.0.1'\\
        description = ''\\
        install\_requires=[]\\[0.75ex]
        \normalfont\scriptsize
        The file is only a package declaration: it exposes no installation-time
        action for code or behavior retrieval to match.
        \vspace{0.5ex}
    \end{minipage}}
    \hfill
    \fcolorbox{black}{black!4}{%
    \begin{minipage}[t]{0.30\textwidth}
        \vspace{0.5ex}
        \textbf{C. Security-sensitive benign tool}\\[-0.25ex]
        \scriptsize\textbf{Package:} \texttt{authcaptureproxy-1.1.4}\\
        \scriptsize\textbf{Gold:} benign \quad \textbf{Outcome:} FP (3/7)\\[0.75ex]
        \rule{\linewidth}{0.3pt}\\[-0.5ex]
        \footnotesize\ttfamily
        name = 'authcaptureproxy'\\
        description = '... capture auth\\
        information from a webpage'\\
        entry\_points = \{'console\_scripts': [\\
        \quad 'auth\_capture\_proxy = ...']\}\\[0.75ex]
        \normalfont\scriptsize
        Benign metadata contains terms that resemble credential-theft signals,
        but the excerpt shows no malicious installation behavior.
        \vspace{0.5ex}
    \end{minipage}}
    \caption{Concrete error examples from the held-out \texttt{setup.py}
    files. FN and FP denote false negative and false positive, respectively;
    the parenthesized count is the number of Llama configurations making that
    error. Excerpts are shortened only for presentation.}
    \label{fig:error-examples}
\end{figure*}

\paragraph{Jev error analysis.} Direct Jev makes 20 false negatives and no
false positives. Behavior guidance corrects six of those direct false
negatives, but it flips eight direct true positives to benign and introduces
one false positive, producing 22 false negatives and one false positive
overall. Figure~\ref{fig:jev-error-examples} shows the three corresponding
failure modes from the saved typed outputs. These probabilities are the model's
decision signals, not calibrated probabilities or evidence of causal model
reasoning.

\begin{figure*}[!ht]
    \centering
    \fcolorbox{black}{black!4}{%
    \begin{minipage}[t]{0.30\textwidth}
        \vspace{0.5ex}
        \textbf{A. Guidance corrects a borderline miss}\\[-0.25ex]
        \scriptsize\textbf{Package:} \texttt{testingijijwdaijdwa}\\
        \scriptsize\textbf{Gold:} malicious \quad \textbf{Direct:} benign (0.52)\\
        \scriptsize\textbf{Guided:} malicious (0.87)\\[0.75ex]
        \rule{\linewidth}{0.3pt}\\[-0.5ex]
        \footnotesize\ttfamily
        install\_requires=[\\
        \quad 'browser\_cookie3', 'discordwebhook',\\
        \quad 'robloxpy', 'requests']\\[0.75ex]
        \normalfont\scriptsize
        Stage one marks \texttt{network\_access} (0.50) and
        \texttt{credential\_access} (0.60) present, reversing a near-tied
        direct decision.
        \vspace{0.5ex}
    \end{minipage}}
    \hfill
    \fcolorbox{black}{black!4}{%
    \begin{minipage}[t]{0.30\textwidth}
        \vspace{0.5ex}
        \textbf{B. Guidance removes a direct detection}\\[-0.25ex]
        \scriptsize\textbf{Package:} \texttt{pyevasive}\\
        \scriptsize\textbf{Gold:} malicious \quad \textbf{Direct:} malicious (0.73)\\
        \scriptsize\textbf{Guided:} benign (0.78)\\[0.75ex]
        \rule{\linewidth}{0.3pt}\\[-0.5ex]
        \footnotesize\ttfamily
        description = 'Python Crypter For Red Teaming'\\
        keywords = ['python', 'crypter',\\
        \quad 'avbypass', 'crypt']\\[0.75ex]
        \normalfont\scriptsize
        No Stage-one behavior exceeds 0.5. The fixed behavior vocabulary thus
        supplies no positive typed indicator despite the direct detection.
        \vspace{0.5ex}
    \end{minipage}}
    \hfill
    \fcolorbox{black}{black!4}{%
    \begin{minipage}[t]{0.30\textwidth}
        \vspace{0.5ex}
        \textbf{C. Guidance creates a false alarm}\\[-0.25ex]
        \scriptsize\textbf{Package:} \texttt{authcaptureproxy-1.1.4}\\
        \scriptsize\textbf{Gold:} benign \quad \textbf{Direct:} benign (0.81)\\
        \scriptsize\textbf{Guided:} malicious (0.51)\\[0.75ex]
        \rule{\linewidth}{0.3pt}\\[-0.5ex]
        \footnotesize\ttfamily
        description = '... capture authentication\\
        information from a webpage'\\
        entry\_points = \{'console\_scripts': [...]\}\\[0.75ex]
        \normalfont\scriptsize
        Stage one marks \texttt{credential\_access} present (0.63), moving a
        direct benign decision just across the malicious threshold.
        \vspace{0.5ex}
    \end{minipage}}
    \caption{Jev decision reversals on held-out \texttt{setup.py} files.
    Parenthesized values are the saved probability of the displayed decision.
    The behavior-guided variant receives the full stage-one vector, not only
    the indicators shown here. Excerpts are shortened for presentation.}
    \label{fig:jev-error-examples}
\end{figure*}

\paragraph{Llama parse failures.} Structured-output decoding fails on 2.3--7.3\%
of samples depending on the method. The no-RAG baseline (6.5\%) and
contrastive method (7.3\%) are worst; both correspond to the largest average
prompt complexity (long obfuscated targets for no-RAG; dual evidence blocks
for contrastive), consistent with backend truncation of long outputs. We
treat coverage as a first-class result rather than silently dropping failed
samples. The Jev comparators use typed outputs instead of JSON generation;
their separate 0.991 coverage is caused by nine unusually large sources that
returned API validation errors, not by output decoding.
\endinput
